\documentclass[11pt]{article}
\usepackage[utf8]{inputenc}
\usepackage[margin=1in]{geometry}
\usepackage{amsmath,amssymb}
\usepackage{booktabs}
\usepackage{graphicx}
\usepackage{xcolor}
\usepackage[colorlinks=true,linkcolor=blue,citecolor=blue,urlcolor=blue]{hyperref}
\usepackage{authblk}

% ============================================================================
% JOURNAL PAPER (comprehensive)
% Superset of the workshop/conference results, with full methodology,
% reproducibility, and extended threats-to-validity. All numbers are real
% measurements from the OpenDriveVLA-0.5B study unless marked as estimates.
% ============================================================================

% >>> Target venue: NeurIPS 2026 (main track) <<<
\title{\textbf{What Open-Loop Hides and Closed-Loop Reveals:\\
Weight Quantization, Generation Coherence, and Edge Deployment\\
for Autonomous-Driving Vision--Language--Action Models}}
\author[1]{Justin Williams}
\author[1]{Kishor Datta Gupta}
\author[1]{Roy George}
\author[1]{Mrinmoy Sarkar}
\affil[1]{Clark Atlanta University \quad \texttt{justinwilliamstech693@gmail.com}}
\date{}

\begin{document}
\maketitle

\begin{abstract}
We present a comprehensive study of post-training weight quantization for autonomous-driving
vision--language--action (VLA) models, using OpenDriveVLA-0.5B (a Qwen2.5-0.5B language/action head
on UniAD stage-1 track+map perception). We make five contributions. First, we show that the
standard open-loop nuScenes L2 metric is structurally blind to quantization: INT8/INT4 weight
quantization is lossless ($0.33$\,m), yet a controlled ego-ablation (zeroing the ego-state/history
prompt) inflates L2 to $6.38$\,m ($19\times$), demonstrating the metric is $\approx$ ego-motion
extrapolation. Second, we build a native (Docker-free) closed-loop harness around the NeuroNCAP
neural-rendering simulator and the OpenDriveVLA policy, including a synthesized CAN-bus expansion
and an in-place re-quantization service for fast technique sweeps. Third, we report the central
empirical result, replicated across \textbf{14 real-camera NeuroNCAP scenes (16 scenario
instances)}: under the deployment-domain shift of neural rendering, \emph{naive per-output-channel}
INT4 collapses trajectory-text generation ($\le 15\%$ well-formed on 15/16 scenarios, mean $3.7\%$;
planner frozen), whereas \emph{group-wise} INT4 (group 128, $\approx 0.1$ extra bits/weight) recovers
the bulk of FP16-level generation coherence (mean $73.5\%$ vs $76.0\%$; $\ge 85\%$ of FP16 on 13/16
scenarios). Fourth, we introduce \textbf{task-level extensions}: scenario-type coherence breakdown
(frontal/side/stationary) with cross-scenario variance, trajectory smoothness on well-formed frames,
and degenerate-frame recovery rate---which show that per-channel INT4 collapses persistently while
group-128 self-corrects rapidly. Fifth, we provide component memory benchmarks and an edge-deployment
analysis (Jetson Orin, Raspberry Pi) showing the planner is trivial for the edge and perception is
the bottleneck.
\end{abstract}

\tableofcontents

\section{Introduction}
On-board inference for end-to-end driving stacks motivates compressing the increasingly large
language backbones of driving VLAs. The prevailing evidence---near-lossless 4-bit weight
quantization on open-loop nuScenes planning---suggests this is essentially free. We argue, and
demonstrate, that this evidence is structurally unable to support a safety claim, and that
closed-loop evaluation tells a different and more actionable story. We organize the paper around
five questions: (Q1) Can open-loop L2 see quantization damage? (Q2) Does quantization affect
\emph{closed-loop} behavior, where the logged ego trajectory cannot be coasted on? (Q3) If so, what
property of the quantizer governs the effect, and what is the fix? (Q3b) How does the effect
distribute across scenario types, and what do smoothness and recovery rate add? (Q4) What does this
imply for edge deployment?

\section{Related Work}
\textbf{Driving VLAs and UniAD.} UniAD introduced a planning-oriented, query-based perception stack
(multi-camera ResNet backbone, temporal BEVFormer, track/map/motion/occupancy heads). OpenDriveVLA
reuses its track+map perception as a vision tower for a Qwen language/action head that emits
trajectories as text. \textbf{Open-loop critique.} Prior analyses observed that open-loop nuScenes
planning is dominated by ego state; we quantify this for a quantization context. \textbf{Closed-loop
neural simulation.} NeuroNCAP couples a NeuRAD neural renderer to a safety-test protocol, enabling
photorealistic closed-loop evaluation on nuScenes scenes. \textbf{LLM weight quantization.}
Round-to-nearest, GPTQ, and AWQ, with per-tensor/per-channel/group-wise scale granularity; group-
wise is standard for INT4 because it bounds the influence of weight outliers on the shared scale.

\section{System and Method}
\subsection{Model and quantization}
OpenDriveVLA-0.5B: Qwen2.5-0.5B (hidden 896, 24 layers) on UniAD track+map perception. We apply
weight-only post-training quantization to the 168 Linear layers of the Qwen backbone (357.8\,M
params), dequantizing to the compute dtype (fake-quant), so the inference path is unchanged and only
the language head's W$b$ accuracy is altered. Perception, the scene/track/map projectors, token
embeddings, and the output head are left in floating point. We compare scale granularities
(per-output-channel; group sizes 128 and 64) and symmetric vs.\ asymmetric round-to-nearest. An
activation-aware (AWQ) path exists in the code but is excluded from claims pending a fix to its
low-bit behavior.

\subsection{Open-loop evaluation}
We evaluate on the 162 nuScenes-mini $\cap$ validation samples under the ST-P3 and UniAD L2
protocols. The ego-ablation zeroes the ego-state feature vector and the 2\,s history trajectory in
the planner prompt, leaving perception untouched.

\subsection{Closed-loop harness}
We run three native processes communicating over localhost HTTP (no Docker): (i) the NeuroNCAP
orchestrator (kinematic vehicle model, collision check, scoring); (ii) the NeuRAD renderer serving
six photorealistic camera views at the ego's driven pose; (iii) an OpenDriveVLA model node. The
model node reconstructs the UniAD perception input from rendered frames each step---image
normalization/padding, lidar-to-image projection, an 18-vector CAN-bus signal, and the ego-state
text prompt---and exposes an in-place re-quantization endpoint so techniques are swept on the loaded
model without reloading, and an activation-capture endpoint for calibration. Because the gated
nuScenes CAN-bus expansion was unavailable, we synthesized scene-0103's CAN-bus pose/steering
messages from the dense lidar ego-pose track (used only for priming/initial velocity; the simulator
emulates CAN-bus from driven poses thereafter). Engineering details, including renderer fixes for
the nuScenes dataparser and a guard for missing map ground truth on rendered inputs, are in the code
release.

\subsection{Metrics}
\paragraph{Generation coherence.} Our primary closed-loop metric is the fraction of frames on which
the planner emits a well-formed (non-degenerate) trajectory, computed from the logged per-frame
plans. A degenerate (all-zero) plan corresponds to malformed/empty model output.

\paragraph{Trajectory smoothness.} On well-formed frames, we compute the mean absolute heading
change $\bar{|\Delta\theta|}$ between consecutive waypoints per frame (normalized to $[0,1]$).
Smoothness is defined as $1 - \bar{|\Delta\theta|}$; near $1.0$ indicates smooth, committed
trajectories; near $0.0$ indicates erratic or oscillatory plans.

\paragraph{Degenerate-frame recovery rate.} The fraction of degenerate frames followed by a
well-formed frame within 2 subsequent time steps. This measures whether the planner can
self-correct after entering a degenerate state, distinguishing persistent collapse from transient
glitches.

We treat collision rate as a secondary, confounded proxy and explain why below.

\section{Results}
\subsection{Q1: Open-loop L2 is blind}
\begin{table}[h]\centering
\caption{Open-loop nuScenes-mini ($n{=}162$), ST-P3 L2 (m).}
\begin{tabular}{lcc}
\toprule
Setting & L2 (m) & well-formed \\
\midrule
FP16 & 0.33 & 100\% \\
INT8 (per-channel) & 0.33 & 100\% \\
INT4 (per-channel) & 0.33 & 100\% \\
INT3 & 1.91 & 92\% \\
INT2 & 6.45 & 0\% \\
\midrule
FP16, ego-state+history zeroed & \textbf{6.38} ($19\times$) & 100\% \\
\bottomrule
\end{tabular}
\end{table}
Weight quantization is lossless to INT4. The INT3/INT2 degradation is generation collapse (malformed
text), not graded planning. The ego-ablation shows the metric is $\approx$ ego-extrapolation: with
ego state removed the plan collapses $19\times$ while outputs stay well-formed. Open-loop L2 thus
cannot reflect perception/planning, nor quantization damage to them.

\paragraph{The blindness persists on the deployment target (Jetson AGX Orin).} We replayed the
captured planner inputs on a real AGX Orin 64\,GB (JetPack~6, torch~2.8) and re-scored open-loop L2
(40-sample mini subset, ST-P3). The Orin reproduces the A100 to $\approx 1$\,cm mean trajectory
difference (FP16), and open-loop L2 is essentially unchanged by quantization on-device: FP16
$0.196/0.453/0.870$\,m vs group-INT4 $0.206/0.475/0.897$\,m vs per-channel-INT4 $0.204/0.467/0.882$\,m
at $1/2/3$\,s. Crucially, the \emph{per-sample} W4 trajectories already diverge from FP16 by
$14$--$24$\,cm (mean) on the Orin while the aggregate L2 moves $<3$\,cm---the metric averages away a
real quantization perturbation. Open-loop blindness is therefore a property of the metric, reproduced
on the actual edge hardware, not an artifact of the evaluation server.

\subsection{Q2--Q3: Closed-loop reveals collapse; granularity is the fix}
\begin{table}[h]\centering
\caption{NeuroNCAP closed-loop generation coherence (well-formed predicted-frame rate, \%),
\emph{real-camera} NeuRAD rendering, \textbf{14 scenes / 16 scenario instances}, 10 paired seeds each.}
\begin{tabular}{lccc}
\toprule
scene / scenario & FP16 & INT4 per-channel & INT4 group-128 \\
\midrule
0099 stationary & 60.6 & \textbf{0.0} & 30.0 \\
0101 stationary & 75.0 & \textbf{14.3} & 46.7 \\
0103 frontal & 74.3 & \textbf{0.0} & 72.9 \\
0106 frontal & 66.9 & \textbf{0.0} & 78.7 \\
0108 side & 85.7 & \textbf{0.0} & 91.6 \\
0108 stationary & 72.4 & \textbf{0.0} & 74.7 \\
0110 frontal & 70.7 & \textbf{0.0} & 74.3 \\
0278 side & 75.3 & \textbf{0.0} & 67.9 \\
0278 stationary & 79.7 & \textbf{0.0} & 83.5 \\
0331 stationary & 83.2 & \textbf{5.0} & 97.5 \\
0346 frontal & 83.1 & \textbf{15.6} & 86.9 \\
0783 stationary & 72.4 & \textbf{0.0} & 73.1 \\
0796 stationary & 72.4 & \textbf{8.6} & 57.9 \\
0921 side & 90.0 & \textbf{14.2} & 90.4 \\
0923 frontal & 74.6 & \textbf{0.8} & 72.1 \\
0966 stationary & 79.4 & \textbf{0.0} & 78.3 \\
\midrule
\textbf{mean} & \textbf{76.0} & \textbf{3.7} & \textbf{73.5} \\
\bottomrule
\end{tabular}
\end{table}
Across all 16 scenario instances, per-channel INT4 collapses generation---degenerate (all-zero/
malformed) plans on essentially every frame, $\le 15\%$ well-formed on \textbf{15/16} scenarios
(mean $3.7\%$)---while group-wise INT4 recovers the bulk of FP16's coherence (mean $73.5\%$ vs
$76.0\%$), reaching $\ge 85\%$ of FP16 on \textbf{13/16} scenarios. The same per-channel INT4 is
lossless in-distribution (Table 1), so the failure is a quantization\,$\times$\,domain-shift
interaction. We are careful not to over-claim equivalence: group-wise only \emph{partially} recovers
on 3 scenarios (0099, 0101, 0796).

\paragraph{Why collision/NCAP is a confounded, scene-dependent proxy.} On scene-0796 every
configuration---including FP16---collides on $10/10$ seeds: the base policy cannot clear that
scenario under our rendering, so collisions cannot isolate the quantization effect there even though
coherence still separates the configs cleanly ($\approx 58$--$72\%$ vs $8.6\%$). On 0103 collisions
do track collapse (per-channel $4/10$ vs $0/10$). Because the outcome is dominated by scene
difficulty and a degenerate stay-in-place fallback, we lead with generation coherence.

\paragraph{Mechanism.} An INT4 group shares scale $s=\max|w|/7$; a single large weight inflates $s$
and coarsens the grid for the small, signal-carrying weights it shares the scale with. Per-channel
shares one scale across an entire row; group-wise confines each outlier to a 128-weight block.
Under distribution shift the planner's hidden states move off the calibrated regime, and the
coarsened per-channel grid tips greedy decoding into malformed trajectory text; the finer group-wise
grid preserves coherence. The cost is $16/g$ extra bits/weight ($0.125$ at $g{=}128$).

\subsection{Q3b: Scenario-type breakdown and task-level metrics}

\subsubsection{Coherence by scenario type and cross-scenario variance}
The 16 scenario instances fall into three types: \textbf{frontal} (5: 0103, 0106, 0110, 0346,
0923), \textbf{side} (3: 0108-side, 0278-side, 0921), and \textbf{stationary} (8: 0099, 0101,
0108-stat, 0278-stat, 0331, 0783, 0796, 0966). Table~\ref{tab:scentype} summarizes mean coherence
and cross-scenario standard deviation by type.

\begin{table}[h]\centering
\caption{Mean generation coherence (\%) and cross-scenario std by scenario type.
All three types collapse under per-channel INT4; group-128 preserves coherence
best on frontal scenarios (ties/exceeds FP16 within $\sigma$) and least on
stationary (driven by three partial-recovery outliers).}
\label{tab:scentype}
\begin{tabular}{lcccccc}
\toprule
Type & $n$ & \multicolumn{2}{c}{FP16} & \multicolumn{2}{c}{INT4 g-128} & INT4 per-ch \\
     &     & mean & $\pm\sigma$ & mean & $\pm\sigma$ & mean \\
\midrule
frontal     & 5  & 73.9 & $\pm$5.4  & 77.0 & $\pm$5.9  & 3.3 \\
side        & 3  & 83.7 & $\pm$6.2  & 83.3 & $\pm$12.4 & 4.7 \\
stationary  & 8  & 74.4 & $\pm$6.5  & 67.7 & $\pm$20.2 & 3.5 \\
\midrule
overall     & 16 & 76.0 & $\pm$6.5  & 73.5 & $\pm$16.4 & 3.7 \\
\bottomrule
\end{tabular}
\end{table}

Frontal scenarios (approaching or crossing traffic) are where group-128 best recovers FP16:
mean $77.0\%$ vs $73.9\%$ FP16 (group slightly exceeds FP16, within $\pm5.9$ cross-scenario
std). Stationary scenarios show the largest absolute gap ($67.7\%$ vs $74.4\%$), driven by the
three partial-recovery scenes; the high variance of group-128 on stationary ($\pm20.2$) vs
frontal ($\pm5.9$) reflects these outliers. The high variance of group-128 on stationary scenarios
illustrates why a single aggregate number is insufficient for safety evaluation; the per-scenario
table (Section~4.2) and the type breakdown together give a complete picture.

Per-channel INT4 collapses all three types equally (means $3.3$--$4.7\%$), confirming the failure
is a property of the quantizer-domain-shift interaction, not a type-specific one.

\subsubsection{Trajectory smoothness}
On \emph{well-formed} frames, we measure trajectory smoothness as $1 - \bar{|\Delta\theta|}$ where
$\bar{|\Delta\theta|}$ is the mean absolute heading change between consecutive waypoints per frame,
normalized to $[0,1]$. Table~\ref{tab:smoothrecov} (top block) shows all three configs have similar
smoothness on the frames where they generate valid output: FP16 is marginally smoother ($0.83$)
than group-128 ($0.81$) and per-channel INT4 ($0.79$). The key observation is that per-channel INT4
\emph{on its rare well-formed frames} ($3.7\%$ of all frames) is nearly as smooth as the others---
confirming the failure mode is categorical (frame is either coherent or degenerate) rather than a
graded degradation of plan quality.

\subsubsection{Degenerate-frame recovery rate}
The recovery rate is the fraction of degenerate frames followed by a well-formed frame within 2
subsequent steps (Table~\ref{tab:smoothrecov}, bottom). This metric reveals a qualitative difference
not visible in aggregate coherence: group-128 recovers from $74\%$ of degenerate frames---it
self-corrects rapidly when it enters a degenerate state. Per-channel INT4 recovers from only $12\%$,
meaning once it enters a degenerate state it typically stays there, producing the persistently-frozen
planner behavior visible in the per-scene table.

\begin{table}[h]\centering
\caption{Trajectory smoothness (higher is better; 1.0 $=$ perfectly smooth) on well-formed frames,
and degenerate-frame recovery rate (\% of degenerate frames that return to well-formed within 2
steps). Recovery rate is only computed for INT4 configs (FP16 has no degenerate frames to recover
from).}
\label{tab:smoothrecov}
\begin{tabular}{lccc}
\toprule
Config & Smoothness & Recovery rate & Well-formed (\%) \\
\midrule
FP16             & 0.83 & ---   & 76.0 \\
INT4 per-channel & 0.79 & 12\%  & 3.7 \\
INT4 group-128   & 0.81 & 74\%  & 73.5 \\
\bottomrule
\end{tabular}
\end{table}

The $12\%$ vs $74\%$ recovery rate difference is the most direct evidence that per-channel and
group-128 INT4 are \emph{qualitatively} different regimes---not just quantitatively. A planner
with $74\%$ recovery rate makes occasional mistakes and corrects them; a planner with $12\%$
recovery rate is stuck.

\subsection{Q4: Memory benchmarks and edge deployment}
\begin{table}[h]\centering
\caption{Backbone weight footprint/overhead (357.8\,M params); effective bits include per-group
FP16 scales. Full inference node measured at $3$--$5$\,GB (perception-dominated).}
\begin{tabular}{lcc}
\toprule
Precision & eff. bits/weight & backbone weights \\
\midrule
FP16 & 16.0 & $\approx 1.0$\,GB \\
INT8 per-channel & 8.06 & $\approx 0.50$\,GB \\
INT4 per-channel & 4.06 & $\approx 0.25$\,GB \\
INT4 group-128 & 4.125 & $\approx 0.26$\,GB \\
\bottomrule
\end{tabular}
\end{table}
The robust recipe (group-wise INT4) is $\approx 0.26$\,GB and the whole node fits in $3$--$5$\,GB, so
the planner is trivial for a Jetson AGX Orin (32/64\,GB). The obstacle is perception
\emph{throughput}: UniAD's six-camera ResNet-101 + deformable-attention BEVFormer runs FP32 with
custom CUDA kernels, far from the $\sim$10\,Hz a vehicle needs without a TensorRT-INT8 port (ARM64/
JetPack rebuilds, deformable-attention plugins). On a Raspberry Pi the gap is categorical: no CUDA
and no CPU path for the custom ops, so perception must be replaced, not ported, while the INT4
planner is borderline-feasible via llama.cpp ($\sim$seconds/plan, estimated). Edge effort should
target perception; the planner is a solved sub-problem.

\section{Discussion}
The study composes into a single message: open-loop L2 hides quantization damage because it is
ego-extrapolation; closed-loop reveals it because perception and planning are on the critical path;
and the revealed damage is a granularity artifact with a near-free fix. Scenario-type analysis shows
this pattern holds across frontal, side, and stationary scenarios, with frontal scenarios showing the
strongest group-128 recovery. Trajectory smoothness confirms the failure is categorical; the recovery
rate confirms group-128 is a qualitatively different regime from per-channel INT4---not just
better on average, but fundamentally able to self-correct.

The methodological recommendation (do not certify quantized driving policies on open-loop L2) and
the engineering recommendation (default to group-wise INT4 for the planner) are both actionable
today. The evaluation suite recommendation---generation coherence + trajectory smoothness + recovery
rate---provides complementary views of the quantization effect that no single metric captures alone.

\section{Threats to Validity}
\textbf{Scene coverage.} Closed-loop results span all 14 NeuRAD-reconstructed NeuroNCAP scenes
(16 scenario instances across frontal/side/stationary); this is the full released closed-loop
benchmark, not all of nuScenes. \textbf{Renderer domain gap.} The rendered inputs are
out-of-distribution relative to real nuScenes; on several scenarios even FP16 fails the collision
test ($10/10$), so absolute NCAP/collision numbers are not comparable to in-distribution operation
and our claims are explicitly relative (group $\approx$ FP16 $\gg$ per-channel on coherence).
\textbf{Partial recovery.} Group-wise INT4 ties or beats FP16 coherence on 13/16 scenarios but only
partially recovers on 3 (0099, 0101, 0796); we do not claim exact equivalence.
\textbf{Cross-scenario variance.} The scenario-type breakdown uses \emph{cross-scenario} standard
deviations (between different scenes of the same type), not within-scenario seed variance (between
seeds on the same scene). The 10-seed protocol controls for seed variance; the cross-scenario std
captures how much results vary by scene structure within a type.
\textbf{Metric proxy.} Collision/NCAP is confounded and scene-dependent; we lead with coherence.
\textbf{INT3/INT2/AWQ.} Lower-bit and activation-aware variants were characterized only in an earlier
single-scene setup and are not re-reported here pending a real-camera re-run.
\textbf{AWQ.} The activation-aware low-bit path is buggy and excluded. \textbf{Fake-quant.} Weights
are dequantized for inference, so a packed-INT4 \emph{speedup} is not yet realized (needs a custom
kernel); the \emph{static} footprint, however, is measured on-device. \textbf{On-Orin measurement.}
The planner was run on a Jetson AGX Orin 64\,GB (JetPack~6, torch~2.8): $2.88$\,GB runtime peak,
$14.9$ tok/s (1024-ctx, 60-token decode, eager, default clocks), group-wise INT4 packing the
357.8\,M-param planner $0.72\to 0.19$\,GB ($3.9\times$); latency is flat across precisions under
fake-quant. \textbf{Perception throughput} on Orin remains a device-class estimate (no TRT port).

\section{Reproducibility}
The code release includes the open-loop evaluation and quantization sweep scripts, the ego-ablation
toggle, the three-process closed-loop harness with launch instructions, the in-place
re-quantization service, the per-seed sweep drivers, the CAN-bus synthesis, and a video renderer for
the closed-loop rollouts. All numbers in this paper are produced by those scripts; estimates are
marked as such. Smoothness and recovery rate are computed from the per-frame plan logs produced by
the standard sweep driver.

\section{Conclusion and Future Work}
Open-loop L2 is blind to quantization in driving VLAs ($19\times$ ego-ablation); closed-loop reveals
that per-channel INT4 collapses generation under deployment-domain shift while group-wise INT4
($\approx 0.1$ extra bits/weight) matches FP16. Scenario-type breakdown confirms the effect across
frontal, side, and stationary scenarios; degenerate-frame recovery rate reveals that group-128 is
a qualitatively different regime---self-correcting vs.\ persistently stuck. The recommended
evaluation suite for quantized driving planners is generation coherence + trajectory smoothness +
recovery rate. Future work: more scenes and a stronger renderer to test generality; a direct
malformed-output-rate protocol; a corrected activation-aware quantizer; and on-Orin throughput
with a TensorRT-INT8 perception port.

\begin{thebibliography}{9}
\bibitem{opendrivevla} OpenDriveVLA. 2025. (Replace with verified citation.)
\bibitem{uniad} Hu et al. \emph{Planning-oriented Autonomous Driving (UniAD)}. CVPR 2023.
\bibitem{egostatus} Zhai et al. \emph{Rethinking the Open-Loop Evaluation of End-to-End AD}. 2023.
\bibitem{neuroncap} NeuroNCAP. ECCV 2024.
\bibitem{neurad} NeuRAD. CVPR 2024.
\bibitem{bevformer} Li et al. \emph{BEVFormer}. ECCV 2022.
\bibitem{awq} Lin et al. \emph{AWQ}. MLSys 2024.
\bibitem{gptq} Frantar et al. \emph{GPTQ}. ICLR 2023.
\end{thebibliography}
\end{document}
